\documentclass[11pt]{article}

\usepackage{amsmath,amssymb,amsthm,mathtools}
\usepackage[margin=1in]{geometry}
\usepackage{hyperref}

\newtheorem{theorem}{Theorem}
\newtheorem{lemma}[theorem]{Lemma}
\newtheorem{definition}[theorem]{Definition}
\newtheorem{remark}[theorem]{Remark}

\DeclareMathOperator{\Hilb}{Hilb}
\DeclareMathOperator{\Inv}{inv}
\DeclareMathOperator{\OP}{OP}
\DeclareMathOperator{\Sym}{Sym}

\newcommand{\C}{\mathbb{C}}
\newcommand{\Sn}{\mathfrak{S}_n}
\newcommand{\qint}[1]{[#1]_q}

\begin{document}
\title{Submission C solution to Problem 9}
\date{}
\maketitle


\begin{abstract}
Let \(R_n^{(m)}\) be the diagonal coinvariant algebra of \(\mathfrak S_n\) in
\(m\) sets of \(n\) commuting variables, and write
\[
\Hilb(R_n^{(m)};q_1,\dots,q_m)
=
\sum_{\lambda} c_\lambda(n)s_\lambda(q_1,\dots,q_m).
\]
For hook shapes \(\lambda=(a+1,1^b)\), we give a positive ordered-set-partition
interpretation of \(c_\lambda(n)\).
\end{abstract}

\section{Statement of the hook formula}

For \(r,k\geq 1\), let \(\OP_{r,k}\) be the set of ordered set partitions
\[
\pi=(B_1\mid B_2\mid \cdots \mid B_k)
\]
of \([r]=\{1,\dots,r\}\) into \(k\) nonempty blocks. Define
\[
\Inv(\pi)
=
\#\{(i,j,x): 1\leq i<j\leq k,\ x\in B_i,\ x>\min(B_j)\}.
\]
Also write \([k]_q=1+q+\cdots+q^{k-1}\).

\begin{theorem}[Hook coefficients]
Let \(n\geq 1\). For \(a,b\geq 0\), write the hook
\[
(a\mid b):=(a+1,1^b).
\]
Then
\[
\boxed{
\sum_{a,b\geq 0} c_{(a\mid b)}(n)\,q^a z^b
=
\sum_{r=1}^{n-1}\sum_{k=1}^{r}
z^{r-k}[k]_q
\sum_{\pi\in \OP_{r,k}} q^{\Inv(\pi)}.
}
\]
Equivalently,
\[
\boxed{
c_{(a+1,1^b)}(n)
=
\#\left\{
(r,\pi,t):
\begin{array}{l}
1\leq r\leq n-1,\ \pi\in \OP_{r,r-b},\\
0\leq t\leq r-b-1,\ \Inv(\pi)+t=a
\end{array}
\right\}.
}
\]
In particular, \(c_{(a+1,1^b)}(n)=0\) unless \(0\leq b\leq n-2\).
\end{theorem}

\section{Ordered set partitions and \(q\)-Stirling numbers}

Let
\[
A_{r,k}(q):=\sum_{\pi\in \OP_{r,k}} q^{\Inv(\pi)}.
\]

\begin{lemma}
The polynomials \(A_{r,k}(q)\) satisfy
\[
A_{r,k}(q)=[k]_q\bigl(A_{r-1,k}(q)+A_{r-1,k-1}(q)\bigr),
\]
with \(A_{1,1}(q)=1\). Hence
\[
A_{r,k}(q)=[k]!_q\,\mathrm{Stir}_q(r,k),
\]
where \(\mathrm{Stir}_q(r,k)\) is the Carlitz \(q\)-Stirling number defined by
\[
\mathrm{Stir}_q(r,k)=[k]_q\,\mathrm{Stir}_q(r-1,k)+\mathrm{Stir}_q(r-1,k-1).
\]
\end{lemma}

\begin{proof}
Remove the largest element \(r\) from an ordered set partition
\(\pi=(B_1|\cdots|B_k)\).

If \(r\) lies in a block \(B_i\) of size at least \(2\), removing \(r\) gives an
ordered set partition in \(\OP_{r-1,k}\). Conversely, inserting \(r\) into block
\(i\) of a partition in \(\OP_{r-1,k}\) creates exactly \(k-i\) new inversions,
since \(r>\min(B_j)\) for every block \(j>i\). Thus these insertions contribute
\[
\sum_{i=1}^k q^{k-i} A_{r-1,k}(q)=[k]_q A_{r-1,k}(q).
\]

If \(r\) is a singleton block, deleting that block gives a partition in
\(\OP_{r-1,k-1}\). Conversely, inserting the singleton block \(\{r\}\) in
position \(i\) creates again \(k-i\) new inversions. These insertions contribute
\([k]_q A_{r-1,k-1}(q)\). Adding the two cases proves the recurrence. The
identification with \([k]!_q\mathrm{Stir}_q(r,k)\) follows from the same recurrence
and initial condition.
\end{proof}

\section{Superspace and hook extraction}

Let \(\Omega_n=\C[x_1,\dots,x_n]\otimes \bigwedge(\theta_1,\dots,\theta_n)\), where
the \(x_i\) commute and the \(\theta_i\) anticommute. Let \(\Sn\) act diagonally by
\[
w\cdot x_i=x_{w(i)},\qquad w\cdot \theta_i=\theta_{w(i)}.
\]
The superspace coinvariant ring is
\[
SR_n:=\Omega_n/\langle(\Omega_n)^{\Sn}_+\rangle.
\]
Let \(q\) track \(x\)-degree and \(z\) track \(\theta\)-degree.

\begin{lemma}[Hook specialization]
Let
\[
C_n(q,z):=\sum_{a,b\geq 0} c_{(a\mid b)}(n)q^a z^b.
\]
Then
\[
\Hilb(SR_n;q,z)=1+(q+z)C_n(q,z).
\]
\end{lemma}

\begin{proof}
The construction of diagonal coinvariants is functorial in the row-space. For a
purely even row-space of dimension \(m\), its character is
\[
\sum_\lambda c_\lambda(n)s_\lambda(q_1,\dots,q_m).
\]
Evaluating the same polynomial functor on a \(1|1\)-dimensional superspace gives
the superspace coinvariant ring \(SR_n\). By the Berele--Regev hook Schur
character theory, the ordinary Hilbert character of the Schur functor
\(\mathbb S_\lambda\) on one even variable \(q\) and one odd variable \(z\) is zero
unless \(\lambda\) is a hook. Moreover,
\[
hs_{(a\mid b)}(q;z)=q^{a+1}z^b+q^a z^{b+1}=(q+z)q^a z^b,
\]
and \(hs_\varnothing(q;z)=1\). Therefore
\[
\Hilb(SR_n;q,z)
=
1+\sum_{a,b\geq 0}c_{(a\mid b)}(n)(q+z)q^a z^b
=
1+(q+z)C_n(q,z).
\]
This is the standard hook Schur specialization; see Berele--Regev
\cite{BereleRegev} or Macdonald \cite[Ch.~I, \S3, Ex.~23]{Macdonald}.
\end{proof}

\section{Hilbert series of \(SR_n\)}

Rhoades and Wilson proved the following Hilbert series formula for the superspace
coinvariant ring \cite[Theorem~1.1 and Eq.~(1.10)]{RhoadesWilson}:
\[
\Hilb(SR_n;q,z)=
\sum_{k=1}^{n} z^{n-k}[k]!_q\,\mathrm{Stir}_q(n,k).
\]
Using the preceding lemma, this can be written as
\[
\Hilb(SR_n;q,z)=
\sum_{k=1}^{n}z^{n-k}A_{n,k}(q).
\]

Set
\[
H_n(q,z):=\sum_{k=1}^{n}z^{n-k}A_{n,k}(q).
\]
Then \(H_1(q,z)=1\).

\begin{lemma}
For \(n\geq 2\),
\[
H_n(q,z)-H_{n-1}(q,z)
=
(q+z)\sum_{k=1}^{n-1} z^{n-1-k}[k]_q A_{n-1,k}(q).
\]
Consequently,
\[
\frac{H_n(q,z)-1}{q+z}
=
\sum_{r=1}^{n-1}\sum_{k=1}^{r}
z^{r-k}[k]_q A_{r,k}(q).
\]
\end{lemma}

\begin{proof}
Using the recurrence from the ordered-set-partition lemma,
\[
A_{n,k}=[k]_q(A_{n-1,k}+A_{n-1,k-1}).
\]
Therefore
\[
\begin{aligned}
H_n
&=
\sum_{k=1}^{n} z^{n-k}[k]_qA_{n-1,k}
+
\sum_{k=1}^{n} z^{n-k}[k]_qA_{n-1,k-1}  \\
&=
\sum_{k=1}^{n-1} z^{n-k}[k]_qA_{n-1,k}
+
\sum_{k=1}^{n-1} z^{n-1-k}[k+1]_qA_{n-1,k} \\
&=
\sum_{k=1}^{n-1} z^{n-1-k}
\bigl(z[k]_q+[k+1]_q\bigr)A_{n-1,k}.
\end{aligned}
\]
Since
\[
z[k]_q+[k+1]_q
=
z[k]_q+1+q[k]_q
=
1+(q+z)[k]_q,
\]
we obtain
\[
H_n
=
H_{n-1}
+
(q+z)\sum_{k=1}^{n-1} z^{n-1-k}[k]_qA_{n-1,k}.
\]
Telescoping from \(H_1=1\) proves the second formula.
\end{proof}

\section{Proof of the theorem}

By the hook-specialization lemma and the Rhoades--Wilson Hilbert series formula,
\[
1+(q+z)C_n(q,z)
=
H_n(q,z).
\]
Thus
\[
C_n(q,z)=\frac{H_n(q,z)-1}{q+z}.
\]
The preceding lemma gives
\[
C_n(q,z)
=
\sum_{r=1}^{n-1}\sum_{k=1}^{r}
z^{r-k}[k]_qA_{r,k}(q)
=
\sum_{r=1}^{n-1}\sum_{k=1}^{r}
z^{r-k}[k]_q
\sum_{\pi\in \OP_{r,k}}q^{\Inv(\pi)}.
\]
Since
\[
[k]_q=\sum_{t=0}^{k-1}q^t,
\]
the coefficient of \(q^a z^b\) counts triples
\[
(r,\pi,t)
\]
such that \(1\leq r\leq n-1\), \(\pi\in \OP_{r,k}\), \(k=r-b\), \(0\leq t\leq k-1\),
and
\[
\Inv(\pi)+t=a.
\]
This is precisely the stated interpretation of \(c_{(a+1,1^b)}(n)\).

\begin{thebibliography}{9}

\bibitem{BereleRegev}
A.~Berele and A.~Regev,
\newblock Hook Young diagrams with applications to combinatorics and to representations of Lie superalgebras,
\newblock \emph{Adv. Math.} \textbf{64} (1987), 118--175.

\bibitem{Bergeron}
F.~Bergeron,
\newblock Multivariate diagonal coinvariant spaces for complex reflection groups,
\newblock \emph{Adv. Math.} \textbf{239} (2013), 97--108.

\bibitem{Macdonald}
I.~G. Macdonald,
\newblock \emph{Symmetric Functions and Hall Polynomials},
\newblock 2nd ed., Oxford University Press, 1995.

\bibitem{RhoadesWilson}
B.~Rhoades and A.~T. Wilson,
\newblock The Hilbert series of the superspace coinvariant ring,
\newblock \emph{Forum Math. Pi} \textbf{12} (2024), Paper No. e16.

\end{thebibliography}

\end{document}